% =====================================================================
%  pacotes.tex  ->  todos os \usepackage do projeto
%  Compilar com pdfLaTeX (codificacao UTF-8) ou, se preferir fontes
%  modernas, com LuaLaTeX/XeLaTeX (basta trocar o bloco de fontes).
% =====================================================================

% ---------- Idioma e codificacao (pdfLaTeX) --------------------------
\usepackage[utf8]{inputenc}
\usepackage[T1]{fontenc}
\usepackage[brazil]{babel}

% ---------- Tipografia -----------------------------------------------
\usepackage{lmodern}      % fonte base limpa; trocar por libertinus se desejar
\usepackage{microtype}    % ajuste fino de espacamento (cara de artigo)

% ---------- Layout da pagina -----------------------------------------
\usepackage[a4paper,margin=2.5cm]{geometry}
\usepackage{setspace}
\usepackage{titlesec}     % titulos de secao estilizados (ver estilo.tex)
\usepackage{enumitem}     % listas mais limpas

% ---------- Cor e caixas didaticas -----------------------------------
\usepackage{xcolor}
\usepackage[most]{tcolorbox}

% ---------- Figuras e diagramas --------------------------------------
\usepackage{graphicx}
\usepackage{float}
\usepackage{subcaption}
\usepackage{tikz}
\usetikzlibrary{trees,positioning,arrows.meta,shapes.geometric,calc,fit,backgrounds}

% ---------- Matematica -----------------------------------------------
\usepackage{amsmath}
\usepackage{amssymb}

% ---------- Simbolos (tesoura da poda, etc.) -------------------------
\usepackage{pifont}

% ---------- Links ----------------------------------------------------
\usepackage[hidelinks]{hyperref}

% ---------- Bibliografia (ABNT) --------------------------------------
% Requer compilacao com biber. Se o Overleaf reclamar do estilo abnt,
% trocar 'style=abnt' por 'style=numeric' temporariamente.
\usepackage[backend=biber,style=abnt]{biblatex}
\addbibresource{referencias/referencias.bib}

% ---------- Pacotes de dominio (ativar quando produzir as figuras) ---
% \usepackage{chessboard}  % tabuleiro das N-Rainhas
% \usepackage{xskak}
% \usepackage{sudoku}      % grade de Sudoku
% \usepackage{animate}     % animacoes embutidas no PDF (Adobe Reader)
